\documentclass[11pt, a4paper]{article}

\usepackage[french]{babel}
\usepackage[T1]{fontenc}
\usepackage[utf8]{inputenc}
\usepackage{array, multirow, tabularx}
\usepackage{graphicx}
\usepackage{multicol}
\usepackage{xcolor}
\usepackage{colortbl}
\usepackage{url}

\usepackage{geometry}
\geometry{
  left=1.5cm,
  right=1.5cm,
  top=2cm,
  bottom=2cm
}

\usepackage{fontspec}
\setmainfont{OpenSans}[Path=./fonts/, Extension=.ttf,
  UprightFont=*-Regular,
  BoldFont=*-Bold,
  ItalicFont=*-Italic]

\usepackage{setspace}
\setstretch{1.2}

\begin{document}

\begin{center}
	\begin{minipage}{0.2\textwidth}
		\includegraphics[height=3cm]{rezoleo_logo.png}
	\end{minipage}
	\begin{minipage}{0.7\textwidth}
		\vspace{0.3cm}
		\Huge \textbf{Charte d'hébergement}\\
		\Huge \textbf{Association \& Projet}
	\end{minipage}
\end{center}

\vspace*{0.5cm}

\hrule
\vspace{.5cm}
\noindent Nom de l'Association/Projet : \dotfill
\vspace{.5cm}\\
Email de Contact : \dotfill
\vspace{.5cm}\\
Nom du compte : \dotfill\\
\hrule
\vspace{1cm}

\noindent Cette charte définit les règles de bonne utilisation de la solution d'hébergement proposée par Rézoléo, et est établie entre l'association/projet indiqué ci-dessus, et l'association Rézoléo (déclarée le 19/05/2017 à la préfecture du Nord sous le RNA W595029406) et représentée par son président.

\section{Respect de la loi}
Les responsables des sites hébergés par le Rézoléo s'engagent à respecter les dispositions légales en vigueur s'imposant à tout usager d'un système informatique, notamment mais sans se limiter à ce qui concerne le RGPD et la loi du 6 janvier 1978 relative à l'informatique, aux fichiers et aux libertés.

\section{Limites à l'utilisation}
Le service d'hébergement du Rézoléo ne peut pas être utilisé pour se livrer à des activités de spam, de publicité commerciale, ou d'activités illégales telles que les attaques \textit{DDOS} ou de \textit{Phishing}. De même, le service d'hébergement du Rézoléo ne peut pas être utilisé pour héberger des fichiers violant la propriété intellectuelle, le droit d'auteur ou le droit à l'image. La mise en place d'un logiciel \textit{serveur}, \textit{relais} ou similaire (p. ex : \textit{bouncer}, \textit{socks}, ...) est formellement interdite sans l'autorisation de Rézoléo.
Le compte dispose d'un espace disque pour ses fichiers. Cette taille d'espace disque est limitée, et peut être changée par Rézoléo. Cette limite est par défaut fixée à \texttt{300Mo}.

\section{Conditions d'accès}
Le droit d'accès à l'espace de stockage est strictement réservé aux membres de l'association/du projet. Ainsi, le mot de passe ne doit être communiqué en aucun cas. Ce compte est incessible et assorti de moyens d'identification. L'accès au compte peut être retiré si les conditions d'accès ne sont plus respectées ou si le comportement de l'utilisateur est contraire à la présente charte. Le responsable du compte doit prévenir Rézoléo de tout accès frauduleux ou tentative d'accès aux ressources qu'il utilise. Il est responsable de la protection de ses fichiers et de l'accès à ses données.

\section{Confidentialité des données}
Les fichiers hébergés sont considérés comme privés. La lecture, la copie ou la modification d'un fichier à des fins autres que de sauvegarde ne peuvent être réalisées qu'après accord du responsable du compte, ou demande des autorités compétentes.

\section{Sanctions et clôture du compte}
L'association Rézoléo est tenue par la loi de signaler toute action illégale constatée. Le non-respect de la présente charte ou de la loi peut entraîner la clôture anticipée du compte, ainsi que toute action légale que Rézoléo jugera nécessaire.
La clôture normale du compte s'effectue soit par demande du responsable du compte, soit 1 an et 6 mois après la dernière signature de la charte ou après la dissolution de l'association ou la fin du projet.

\section{Notifications de changement}
Le responsable du compte s'engage à notifier Rézoléo de tout changement de responsable, qui est susceptible de nécessiter une nouvelle signature de cette charte. De même, en cas de changement organisationnel modifiant la liste des personnes susceptibles d'acéder à l'hébergement, le responsable du compte s'engage à notifier Rézoléo.

\vspace{.5cm}

\noindent Chaque signataire reconnaît en signant avoir pris conscience du contenu de cette charte, des sanctions encourues, et s'engage à la respecter.

\begin{table}[h]
	\centering
	\renewcommand{\arraystretch}{2}
	\small{
		\begin{tabular}{|p{1.5cm}|p{7cm}|p{7cm}|}
			\hline
			\rowcolor{black!10} Promo & NOM Prénom & Date et Signature \\
			\hline
			                          &            &                   \\
			\hline
			                          &            &                   \\
			\hline
			                          &            &                   \\
			\hline
			                          &            &                   \\
			\hline
			                          &            &                   \\
			\hline
			                          &            &                   \\
			\hline
			                          &            &                   \\
			\hline
		\end{tabular}
	}
\end{table}

\end{document}
